\section{Related Work}
\label{sec:related}

\paragraph{Traffic accident anticipation.}
The accident anticipation literature has focused primarily on improving
prediction accuracy from video. The seminal DSA-RNN~\citep{ref4} established
the DAD benchmark and showed that soft attention over detected objects enables
anticipation roughly 2\,s before impact. Subsequent work explored richer
interaction modelling, better spatio-temporal fusion, and larger-capacity
architectures. Examples include agent-centric risk assessment and risky-region
localisation~\citep{ref45}, uncertainty-aware relational reasoning~\citep{ref1},
adaptive learning for rare incidents~\citep{ref1}, dynamic spatio-temporal
attention~\citep{ref17}, graph-based interaction modelling~\citep{ref32,ref36,ref37},
and recent high-capacity fusion pipelines such as CRASH~\citep{liao2024crash},
CCAF-Net~\citep{liu2025ccaf}, and LATTE~\citep{ZHANG2025103173}. Recent work
has also begun to target timing more explicitly through temporal occurrence
prediction, risk propagation, and other formulations that model warning timing
more directly than a pure per-frame binary score. Other recent directions bring
language or textual supervision into anticipation, which is relevant to our use
of concept pseudo-labels even when the mechanism is not an explicit concept
bottleneck.

Despite their architectural diversity, most existing methods still optimize
anticipation mainly through prediction quality, with timing handled indirectly
through the shape of the confidence trajectory rather than through an explicit,
interpretable decision state. This has produced impressive AP numbers, but it
leaves two safety-critical gaps unresolved: the models are rewarded mainly for
being correct rather than early in a decision-theoretic sense, and the
internal evidence driving the alert remains largely latent.

\paragraph{Interpretability in accident anticipation.}
Table~\ref{tab:interp_comparison} (Experiments) makes the key point clear:
existing approaches, when they provide explanation at all, do so in ways that
remain external to the actual decision pathway. \citet{ref2} (DRIVE, MM'21) is
the closest prior work because it couples anticipation with reinforcement
learning and visual explanation. However, its Grad-CAM maps are still
\emph{post-hoc}: they are generated after prediction and do not constitute a
semantic state through which the decision is made. \citet{liao2024and}
(W3AL, MM'24) uses an LLM to verbalise accident location and timing, which is
interesting for description but fundamentally different from exposing a
human-auditable risk representation that directly controls the warning policy.
Other explanation-oriented studies remain similarly limited to heatmaps,
localisation, or auxiliary user-facing outputs~\citep{ref102,karim2023attention}.

This distinction matters. In a safety setting, an explanation is most useful
when it is not merely attached to a decision but structurally involved in the
decision itself. \method{} is designed around exactly this requirement: the
concept activations are not generated after the alert, but are part of the
forward path that produces it. This is why we emphasise \emph{intrinsic}
WHY/WHEN interpretability rather than post-hoc explanation quality alone.

\paragraph{Actor-Critic for temporal decisions.}
\citet{sutton2018reinforcement} establish the Actor-Critic framework.
\citet{ref2} (DRIVE) applies REINFORCE to learn a spatial attention
policy for object selection, using RL as a training signal rather than
for optimizing the final alert rule. Recent timing-oriented accident
anticipation methods have also explored temporal occurrence prediction and
related formulations without exposing an explicit semantic decision state. The
main distinction of \caac{} in the current paper is therefore not simply that
it uses RL, but that the state $\mathbf{s}_t=[\mathbf{h}_t\|\mathbf{c}_t]$
contains concept activations and thus provides a direct route from semantic
risk evidence to timing control.

\paragraph{Concept Bottleneck Models and knowledge distillation.}
Concept Bottleneck Models (CBMs) constrain predictions through interpretable
bottlenecks~\citep{ref102}, offering human-readable intermediate representations
that support user intervention.
Applying CBMs to \emph{video} accident anticipation is non-trivial:
concepts must be dynamic across frames, and temporal dependencies must be
preserved through the bottleneck.
\method{} addresses both challenges through \cgta{} (linking temporal
attention to concept dynamics) and the \caac{} state design (placing
concepts directly inside the MDP state).

For concept supervision, \method{} leverages CLIP~\citep{radford2021learning},
which provides rich semantic visual representations from natural language
supervision and enables pseudo-labelling of 837 traffic-relevant concepts
without manual annotation. This choice should be read as a practical semantic
bootstrap rather than as a claim that the resulting pseudo-concepts are fully
noise-free; in the current paper, the strongest support for the concept layer
comes from architectural transparency, ontology analysis, and downstream
ablation evidence rather than from a standalone concept-label benchmark.
The \tsd{} component distils CLIP future-frame features into current
representations, providing anticipatory semantic grounding inspired by
self-predictive representation learning~\citep{fang2022traffic}.
